\section{Conclusions}\label{sec:conc}
%
Lemma~\ref{lem:solveCUG} shows that a hypothetical polynomial-time algorithm for computing $Z(G;w)$ exactly when 
\[
\log w = \frac{2}{1-\eps-\delta}\frac{\log k}{\Delta}
\]
would refute the UGC. 
This problem is likely \#P hard so we resort to approximation, which we can only do for some range of $w$.
In Theorem~\ref{thm:zero-free} we have $\log w = (1-o(1))\log(k)/\Delta$ when $k$ is small enough that $\log k=o(\Delta)$, and in Theorem~\ref{thm:alg,cluster} we have $\log w = (2-o(1))\log(k)/\Delta$ when $k$ is larger than some $\Delta^{\poly(\Delta)}$. 
This means we must increase $f$ from $k^{-1}$ to solve any CUG problems, and the size of $w$ in these results is what gives the bound on $f$ in Theorems~\ref{thm:main,interp} and~\ref{thm:main,cluster}.
It is therefore important to determine the threshold $f^*$ such that \CUG[f] is equivalent to UG when $f\le f^*$. 
Trivially we have $k^{-1}\le f^*\le 1$, but if one could show e.g.\ that $f^*\ge k^{2\theta-1}$ then Theorem~\ref{thm:main,cluster} would mean the Unique Games problem is in P for bounded-degree graphs when $k$ is large enough and $\eps+\delta < \theta$.
%
%
\subsection{Phase transitions}
%
In this subsection we focus on the ferromagnetic Potts model, which is the special case of our partition function $Z(G;w)$ when the constraints on each edge are the identity permutation and we take $w\ge 1$.
The behaviour of the Potts model on bounded-degree graphs is strongly related to the \emph{phases} of the model on the infinite $\Delta$-regular tree $\TT[\Delta]$. 
We will not define precisely what we mean by a phase or a phase transition here, but as the parameter $w$ varies, the behaviour of the model undergoes certain changes that seem to affect both zeros of the partition function and the dynamics of associated Markov chains.
Häggström~\cite{Haggstrom96} showed that the \emph{uniqueness phase transition} on $\TT[\Delta]$ occurs at $w=w_u(k,\Delta)$, the unique value of $w$ for which the polynomial 
\begin{equation}\label{eq:def,wu}
(k-1)x^\Delta + (2-w-k)x^{\Delta-1}+wx-1
\end{equation}
has a double root in $(0,1)$. 
There is a further \emph{ordered/disordered phase transition} (see~\cite{GSVY16}) at 
\[
w_o(k,\Delta):=\frac{k-2}{(k-1)^{1-2/\Delta}-1}\,.
\]
Below we relate the values of $w$ found in Theorems~\ref{thm:alg,cluster} and~\ref{thm:zero-free} to these phase transitions.
%
\subsection{Potential barriers to improving Theorem~\ref{thm:main,interp}}\label{sec:barrers,interp}
%
To strengthen Theorem~\ref{thm:main,interp} to a result that would refute the UGC, we need roughly a factor two improvement in the leading constant in $\log w^* \sim \log (k)/\Delta$ as $k\to\infty$ (provided $\log k=o(\Delta)$) from Theorem~\ref{thm:zero-free}, but there are reasons to believe it may be hard to make such an improvement. 
%
The authors of~\cite{BGP16} analysed a natural Markov chain known as the \emph{Glauber dynamics} which walks the set of possible colorings of a graph $G$, and when this mixes rapidly we expect an efficient, randomised approximation algorithm for the Potts model partition function to follow.
They showed for the Potts model that Glauber dynamics mixes rapidly on graphs of maximum degree $\Delta$ when 
\[
\log w \le (1+o(1))\frac{\log k}{\Delta-1}\,,
\]
as $k\to\infty$, and that on almost all $\Delta$-regular graphs (for $\Delta\ge 3$), Glauber dynamics mixes slowly when $w$ is just a little larger, satisfying
\[
\log w > (1+o(1))\frac{\log k}{\Delta-1-\frac{1}{\Delta-1}}\,.
\]
These bounds sandwich the phase transition point $w_u$; they also showed that as $k\to\infty$, 
\[
\log w_u = \frac{\log k}{\Delta-1} + O(1)\,.
\]
Thus it appears that $w_u$ is a barrier for approximating $Z(G;w)$ via Glauber dynamics. 
%
Similarly, we expect that $w_u$ is a barrier for the zero-free region. 
Underpinning Lemma~\ref{lem:kcol} is a complex dynamical system (see~\cite{PR18} for a treatment of the case $k=2$), and equation~\eqref{eq:def,wu} appears in the analysis of this system.
Essentially, the behaviour of a fixed point in the complex dynamics changes at $w=w_u$ in a way which means it is reasonable to expect zeros of $Z(G;w)$ to accumulate near $w_u$ for some $G$ with maximum degree $\Delta$. 
Thus we suspect that the method cannot work for $w>w_u$. 
%
%
\subsection{Potential barriers to improving Theorem~\ref{thm:main,cluster}}\label{sec:barrers,cluster}
%
There are several regimes of interest for the algorithm in Theorem~\ref{thm:alg,cluster} that gives Theorem~\ref{thm:main,cluster}.
When we apply Theorem~\ref{thm:alg,cluster} to solve CUG problems, we only need an approximation with $\alpha = Cn$ for constant $C$ in which case the running time is bounded by $kn^{O(1)}e^{O(\Delta)}$. 
Recall that we also need $k\ge \Delta^{O(\Delta^{3/2}})$, and hence when $\Delta$ and $k$ do not grow too fast with $n$, the running time is sub-exponential in $n$. 
For the hypercube with $\Delta=\log n$ and with $k$ as small as the result allows, the algorithm is quasi-polynomial.
In the case where $\Delta$ is constant, to get an approximation as accurate as $\alpha$ being constant the running time of the algorithm is polynomial in $n$.
It is interesting to compare this with a paper of Galanis et al.~\cite{GSVY16} who show that it is \#BIS-hard to approximate the partition function of the Potts model with $k$ colours on graphs of maximum degree $\Delta$ when $w>w_o$. 
With $\zeta=8/\sqrt\Delta$, if we take $k=k_0=\exp((18\Delta+4\Delta\log\Delta)/\zeta)$, then Theorem~\ref{thm:alg,cluster} shows that we can approximate the Potts model partition function on graphs of maximum degree at most $\Delta$ with $w=k_0^{(2-\zeta)/\Delta}$. 
A quick calculation shows that, as $\Delta\to \infty$ (and hence $k\to\infty$), $w_o(k_0,\Delta)\sim k_0^{2/\Delta}$. 
We conclude that if one assumes that there are no efficient algorithms for approximating \#BIS-hard problems, Theorem~\ref{thm:alg,cluster} is very close to optimal in this regime.
%
%
